\section{Discussion and Limitations}
\label{sec:discussion}

\subsection{Stable storage is the reusable interface}

LatchMoE's main lesson is that dynamic object identity and graph replay are not
inherently incompatible. The conflict arises when identity changes are exposed
as changes to graph-visible allocations or execution structure. A fixed
storage interface moves that dynamism into contents and metadata that can be
updated at a controlled boundary. The same principle may apply to other
captured operators with a bounded, replaceable working set, provided that the
runtime can identify a legal staging boundary and enforce publication and
reuse ordering. Our experiments establish this lesson only for the implemented
MoE path.

The versioned lease is equally important as the fixed allocation. Stable
pointers alone do not prevent a transfer completion for an old owner from
publishing stale contents, nor do they prevent eviction during in-flight
compute. Address stability supplies replay compatibility; owner/version checks
and stream dependencies supply temporal correctness. Treating these as one
contract makes failures observable and testable.

\subsection{Control policy is outside the data plane}

LatchMoE accepts an ordered placement plan but constrains it to the actual
routed expert set. This boundary lets future caching or prediction policies
reuse the fixed slots, transfer engine, wave executor, and graph checks. The
current evaluation does not compare placement policies and should not be read
as evidence that the default ordering is optimal. Such a study requires a
shared trace, identical storage capacity, and policy-specific hit, transfer,
and latency measurements on the same data plane.

\subsection{Scope of the current evidence}

The online evidence covers one unquantized model, one Ascend 910B2, tensor
parallelism one, a 14-GiB offload target, 32 slots, \texttt{max\_num\_seqs=1},
client concurrency one, and disabled prefix caching. It validates exactness,
graph replay, lifecycle ordering, bounded residency, and service release in
that configuration. The matched result further shows lower TTFT for multi-wave
than full-layer staging under the same narrow contract. It does not show a
general TPOT improvement.

Several dimensions remain open. Higher concurrency changes the union of active
experts and may alter both slot pressure and transfer--compute overlap. A
different model can change expert shapes, routing distributions, and the
number of managed layers. Other slot counts and memory budgets can change wave
count and KV-cache headroom. Prefix-cache reuse is disabled because the current
correctness contract does not cover cache hits across the offload staging
boundary. Multi-device expert parallelism introduces device-to-device
communication and is not represented by the single-device host-to-device
path.

Before making a broad serving claim, the evaluation must compare full
residency, native prefetch offload, the legacy layered path, and LatchMoE under
the same graph-only manifests. It must also add model, capacity, workload, and
concurrency sensitivity, retain graph or capacity failures as results, and
report dispersion across independent starts. Until that matrix is complete,
the paper's evidence supports a graph-compatible mechanism and the scoped
multi-wave TTFT result.

